\uuid{X5Gb}
\titre{Trigonométrie}
\auteur{bodin ; hanani ; mzari}
\organisation{exo7}
\licence{CC BY-NC-SA 4.0}
\datecreate{2020-01-20}
\niveau{L1}
\module{Analyse}
\chapitre{Fonctions circulaires et hyperboliques inverses}
\sousChapitre{Autre}
\difficulte{1}
\theme{Fonctions usuelles, trigonométrie}
\exo7qcmid{lille-1.fonctions-usuelles.006}

\contenu{
\begin{question}
\qclassification{126.00}
\qtags{motcle=trigonométrie}
Soit $f(x)=  \arcsin (\cos  x)$ et  $ g(x)= \arccos (\sin x) $. Quelles sont les assertions vraies ?
\begin{answers}
    \bad{$f$ est périodique de période $\pi$.}

    \good{$g$ est périodique de période $2\pi$.}

    \good{$f$ est une fonction paire.}

    \bad{$g$ est une fonction impaire.}
\end{answers}
\begin{explanations}
Les fonctions $x\ \mapsto   \sin x$ et $x  \mapsto  \cos  x$ sont périodiques de période $2\pi$, donc $f$ et $g$ sont de période $2\pi$. La fonction $x  \mapsto  \cos  x$  est paire, donc $f$ l'est aussi. La fonction $x  \mapsto  \sin   x$ est impaire, mais la fonction $x  \mapsto   \arccos  x$ n'est ni paire ni impaire, donc $g$ n'est ni paire ni impaire.
\end{explanations}
\end{question}
}
